\chapter{Algebraic structure: symmetry, brackets, and roots}
\label{learn:algebra}

\begin{learningcard}{Prerequisites and destination}
Use proof, vector spaces, matrices, and eigenvectors from
Chapters~\ref{learn:proof} and~\ref{learn:linear_algebra}. The aim is to
recognize what an algebraic classification classifies before attaching a
physical interpretation to it. We move from rotations of a triangle to a
small root system, then identify concrete features of the exceptional
$E_8$ endpoint.
\end{learningcard}

\section{A group records reversible composition}
A group is a set $G$ with an associative operation, an identity element,
and an inverse for every element; the operation stays in $G$. If the
operation also commutes, the group is abelian. The rotations of an
equilateral triangle by multiples of $120$ degrees form a three-element
group. Label them $0,1,2$ and compose by addition modulo three: two
successive $120$-degree turns give label two, and three give the identity.
Rotations by arbitrary angles about one fixed axis similarly compose by
angle addition modulo a full turn.

A permutation is a bijection from a set to itself. For three labeled
vertices, the six permutations form the symmetric group $S_3$ under
composition. Write $(12)$ for swapping labels one and two. With the rightmost
permutation acting first, $(12)(23)$ sends one to two, two to three, and
three to one. Reversing the order gives the inverse cycle. Therefore $S_3$
is nonabelian even though its subgroup of three rotations is abelian.

A group action on a set assigns each group element a transformation so that
identity and composition agree. Specifying an action matters: naming a group
beside a physical system does not define how that group's elements transform
the system's states or observations.

\section{Structures supply different operations}
A field has addition and multiplication satisfying the familiar distributive
and associative laws, commutative multiplication, distinct zero and one, and
a multiplicative inverse for every nonzero element. The real and complex
numbers are fields. A vector space over a field has vector addition and
scalar multiplication; it need not multiply its vectors together. An algebra
over a field is a vector space with an additional bilinear product, meaning
a product linear in each argument separately. Extra adjectives specify
whether that product is associative, commutative, or has an identity.

Real square matrices form an associative algebra under multiplication and a
vector space under addition, but they form a multiplicative group only after
restricting to invertible matrices. All matrices together fail that group
condition because singular matrices have no inverse. Keeping these operations
separate prevents a dimension count from substituting for a structural proof.

\section{Worked example: a bracket measures failure to commute}
For matrices define $[A,B]=AB-BA$. This bracket is bilinear,
$[A,A]=0$, and satisfies the Jacobi identity
\[
 [A,[B,C]]+[B,[C,A]]+[C,[A,B]]=0.
\]
To verify Jacobi, expand each nested bracket: the resulting twelve terms
cancel in pairs using associativity. A vector space with a bilinear bracket
satisfying these identities is a Lie algebra.

Consider the trace-zero matrices
\[
 E=\begin{pmatrix}0&1\\0&0\end{pmatrix},\quad
 F=\begin{pmatrix}0&0\\1&0\end{pmatrix},\quad
 H=\begin{pmatrix}1&0\\0&-1\end{pmatrix}.
\]
Multiplication gives $EF=\operatorname{diag}(1,0)$ and
$FE=\operatorname{diag}(0,1)$, hence $[E,F]=H$. Computing the other
products gives $[H,E]=2E$ and $[H,F]=-2F$. These three matrices span
the trace-zero two-by-two matrices, conventionally denoted
$\mathfrak{sl}_2$ over the chosen real or complex field. The eigenvector
language from linear algebra reappears: under the linear map
$X\mapsto[H,X]$, the vectors $E,F$ have eigenvalues two and minus two.

\begin{learningcard}{Historical situation: classification requires definitions}
Cartan's 1894 dissertation investigates the structure of finite continuous
transformation groups. The historical object is a published structural
investigation, not a claim that every later application was contained in that
work. The dissertation belongs to the development from continuous
transformations toward the Lie-algebra classifications used in the advanced
review \citep{cartan1894}. Our matrix brackets provide a small calculation
through which a reader can see what such structural questions ask.
\end{learningcard}

\section{Worked example: the six roots of \texorpdfstring{$A_2$}{A2}}
Let $e_1,e_2,e_3$ be the standard unit vectors in $\RR^3$. The six vectors
$e_i-e_j$ with $i\ne j$ lie in the plane $x_1+x_2+x_3=0$.
They form the root system called $A_2$. Choose simple roots
$\alpha_1=e_1-e_2$ and $\alpha_2=e_2-e_3$; the remaining positive root
is $\alpha_1+\alpha_2=e_1-e_3$, with negatives supplying the other three.
Both simple roots have squared length two and their dot product is minus one.
The angle between them therefore has cosine $-1/2$.

\begin{figure}[htbp]
  \centering
  \begin{tikzpicture}[x=1.65cm,y=1.65cm,>=Stealth,font=\small,line width=0.8pt]
  \coordinate (first) at ({sqrt(2)},0);
  \coordinate (sum) at ({sqrt(2)/2},{sqrt(6)/2});
  \coordinate (second) at ({-sqrt(2)/2},{sqrt(6)/2});
  \coordinate (minusfirst) at ({-sqrt(2)},0);
  \coordinate (minussum) at ({-sqrt(2)/2},{-sqrt(6)/2});
  \coordinate (minussecond) at ({sqrt(2)/2},{-sqrt(6)/2});
  \draw[dashed] (first)--(sum)--(second)--(minusfirst)--(minussum)--(minussecond)--cycle;
  \foreach \point in {first,sum,second,minusfirst,minussum,minussecond} {
    \draw[->,line width=1.1pt] (0,0)--(\point);
  }
  \node[right] at (first) {$\alpha_1$};
  \node[above right] at (sum) {$\alpha_1+\alpha_2$};
  \node[above left] at (second) {$\alpha_2$};
  \node[left] at (minusfirst) {$-\alpha_1$};
  \node[below left] at (minussum) {$-(\alpha_1+\alpha_2)$};
  \node[below right] at (minussecond) {$-\alpha_2$};
  \draw[->] (0.45,0) arc[start angle=0,end angle=120,radius=0.45];
  \node[fill=white,inner sep=1pt] at (0.54,0.3) {$120^\circ$};
  \node[below] at (0,-0.1) {$0$};
  \node[align=center] at (0,-2)
    {All six arrows have length $\sqrt2$.\\
     $\alpha_1\cdot\alpha_2=-1$; $\cos120^\circ=-1/2$.};
\end{tikzpicture}

  \caption{Read the two simple-root arrows first: their angle is 120 degrees. Their vector sum supplies the upper-right root, and negation supplies the opposite arrows. The picture uses orthonormal coordinates in the root plane; the simple roots themselves are not orthogonal.}
  \label{fig:learn:root_hexagon}
\end{figure}

For nonzero roots, reflection perpendicular to $\alpha$ is
$s_\alpha(v)=v-2(v\cdot\alpha)\alpha/(\alpha\cdot\alpha)$.
Reflection in $\alpha_1=e_1-e_2$ swaps the first two coordinates, so it
permutes the six-root set. The Cartan matrix records relative inner products,
$C_{ij}=2(\alpha_i\cdot\alpha_j)/(\alpha_j\cdot\alpha_j)$; here
\[
 C=\begin{pmatrix}2&-1\\-1&2\end{pmatrix},\qquad\det C=3.
\]
The root lattice $Q$ consists of integer combinations of the simple roots.
The weight lattice $P$ consists of vectors pairing integrally with each
coroot $2\alpha_j/(\alpha_j\cdot\alpha_j)$. Solving those pairings
introduces thirds because
\[
 C^{-1}=\frac13\begin{pmatrix}2&1\\1&2\end{pmatrix}.
\]
For this full-rank finite root system, $P/Q$ has three elements. It is a
lattice quotient rather than an automatic selection rule for a fluid.

\section{A concrete endpoint: the \texorpdfstring{$E_8$}{E8} count and lengths}
A standard coordinate construction in $\RR^8$ has two kinds of roots:
\[
 (\pm1,\pm1,0,0,0,0,0,0)\text{ with all position choices},
 \qquad \tfrac12(\pm1,\ldots,\pm1)\text{ with even minus parity}.
\]
There are $\binom82\cdot4=112$ roots of the first kind. Of the $2^8$
sign patterns of the second kind, exactly half have even minus parity,
giving 128. The total is 240. Every first-kind root has squared length
$1+1=2$; every second-kind root has squared length $8/4=2$.
The two families therefore share the same length $\sqrt2$.

The $E_8$ Cartan diagram is a tree with one central vertex and three arms
of lengths one, two, and four. Its matrix has diagonal entries two and
entries minus one along edges. An arm of length $m$ has determinant $m+1$;
this follows from the tridiagonal recurrence $D_m=2D_{m-1}-D_{m-2}$,
$D_0=1,D_1=2$. The inverse entry at its attaching end is $m/(m+1)$,
by the cofactor formula. Eliminating the arms thus gives
\[
 \det C=2\cdot3\cdot5
 \left(2-\frac12-\frac23-\frac45\right)=1.
\]
The corresponding root and weight lattices coincide. Classification identifies
this root system with the rank-eight complex simple Lie algebra $E_8$,
whose dimension is eight Cartan directions plus 240 one-dimensional root
spaces, totaling 248. Proving that classification and constructing all
representations belong to further Lie theory. The finite calculations here
verify the stated count, lengths, and determinant, rather than that entire
classification theorem \citep{kac1990infinite}.

\section{Exercises with full solutions}
\paragraph{1. Compute an inverse in the triangle rotation group.}
Find the inverse of label two under addition modulo three.

\paragraph{Solution.}
An inverse label $a$ satisfies $2+a=0$ modulo three. The label one works
because $2+1=3$, congruent to zero. Geometrically a $120$-degree turn
undoes a $240$-degree turn up to one full revolution.

\paragraph{2. Check a bracket combination.}
Compute $[H,E+F]$ for the matrices above.

\paragraph{Solution.}
Bilinearity and the computed relations give
$[H,E+F]=[H,E]+[H,F]=2E-2F$.
Thus $E+F$ is not an eigenvector of this map: its two components acquire
opposite factors, and $2E-2F$ is not a scalar multiple of $E+F$.

\paragraph{3. Reflect an $A_2$ root.}
Find $s_{\alpha_1}(\alpha_2)$ and check its squared length.

\paragraph{Solution.}
Because $\alpha_2\cdot\alpha_1=-1$ and $\alpha_1^2=2$,
$s_{\alpha_1}(\alpha_2)=\alpha_2+\alpha_1=e_1-e_3$.
Its squared length is two, so the result is another root of the same length.
Chapter~\ref{learn:advanced} extends the distinction between algebraic
structure and physical interpretation.
